
\section{Introduction}
\vspace{-12pt}
A systematic review is a focused literature review that synthesises all relevant literature for a specific research topic. Identifying relevant publications for medical systematic reviews is a highly tedious and costly exercise, often involving multiple reviewers to screen (i.e., assess) upwards of tens of thousands of studies. It is a standard practice to screen each study retrieved for a systematic review by a Boolean query. However, in recent years, there has been a dramatic rise in Information Retrieval methods that attempt to re-rank this set of studies for a variety of reasons, such as stopping the screening early (once a sufficient number of studies have been found) or beginning downstream phases of the systematic review process earlier (such as the acquisition of the full-text of studies). However, a known problem with many of these methods is that they use a different query from the Boolean query used to perform the initial literature search. Instead, most methods typically resort to less informationally representative sources for queries that can be used for ranking, such as the title of the systematic review, e.g.,~\cite{miwa2014reducing} (containing narrow information about the retrieval topic), or concatenating the clauses of the Boolean query together, e.g.,~\cite{alharbi2017ranking} (negating the structural information in Boolean clauses). We instead turn our attention to methods that use more informative sources of information to perform re-ranking. 

Indeed, we focus this reproducibility study on one such method: seed-driven document ranking (SDR) from Lee and Sun~\cite{lee2018seed}. SDR exploits studies that are known a priori to develop the research focus and search strategy for the systematic review. These studies are often referred to as `seed studies' and are commonplace in the initial phases of the systematic review creation process. This method and others such as CLF~\cite{scells2020clf} (which directly uses the Boolean query for ranking) have been shown to significantly outperform other methods that use a na\"ive query representation. Despite this, the SDR method was published when there was little data for those seeking to research this topic, and there have been methods published since that did not include SDR as a comparison. To this end, we devise the following research questions (RQs) to guide our investigation into why we are interested in reproducing the SDR method:

% limitations: one dataset, single seed study


%\subsection{Research Questions}
\vspace{-8pt}
\begin{description}[leftmargin=4pt]
\item[RQ1] \textit{Does the effectiveness of SDR generalise beyond the CLEF TAR 2017 dataset?} The original study was only able to be investigated on a single dataset of systematic review topics. In this study, we plan to use our replicated implementation of SDR to examine the effectiveness of this method across more recent datasets, and datasets that are more topically varied (CLEF TAR 2017 only contains systematic reviews about diagnostic test accuracy).
%While these are indeed a challenging kind of systematic review to produce and of great importance to develop methods to assist with their creation, there are many other kinds of systematic reviews of varying difficulty to produce).
\item[RQ2] \textit{What is the impact of using multiple seed studies collectively on the effectiveness of SDR?} The original study focused on two aspects of their method: an initial ranking using a single seed study and an iterative ranking which further uses the remaining seed studies one at a time. We focus on investigating the first aspect concerning the impact of multiple seed studies (multi-SDR) used collectively for input to produce an initial ranking.
%The original study only utilised a single seed study to produce a ranking.
\item[RQ3] \textit{To what extent do seed studies impact the ranking stability of single- and multi-SDR?} 
%Finally, we examine the impact of seed studies on the variability of SDR effectiveness. 
In a recent study by Scells et al.~\cite{scells2020comparison} to generate Boolean queries from seed studies, it was found that seed studies can have a considerable and significant effect on the effectiveness of resulting queries. We perform a similar study that aims to measure the variance in effectiveness of SDR in single- and multi- seed study settings.
\end{description}

%\noindent
%With the investigation into the above research questions, we will:
%\begin{itemize}[leftmargin=*]
%\item Demonstrate the \textbf{novelty} of the method by performing experiments on more datasets (\textbf{RQ1}), and experiments that reveal more about the effectiveness of the method (\textbf{RQ2}, \textbf{RQ3}).
%\item Assess the \textbf{impact} of SDR towards the Information Retrieval community and the wider systematic review community.
%\item Investigate the \textbf{reliability} of SDR by comparing it to several baselines, including those by participants at the CLEF TAR shared tasks, on publicly available datasets.
%\item Make our complete reproduced implementation of SDR publicly \textbf{available} for others to use as a baseline in future work on re-ranking for systematic reviews.
%\end{itemize}
\noindent
With the investigation into the above research questions, we will (1) demonstrate the \textbf{novelty} of the method by performing experiments on more datasets (\textbf{RQ1}), and experiments that reveal more about the effectiveness of the method (\textbf{RQ2}, \textbf{RQ3}), (2) assess the \textbf{impact} of SDR towards the Information Retrieval community and the wider systematic review community, (3) investigate the \textbf{reliability} of SDR by comparing it to several baselines on publicly available datasets, and (4) make our complete reproduced implementation of SDR publicly \textbf{available} for others to use as a baseline in future work on re-ranking for systematic reviews.
